% ECONOMICS_PROBLEM: {"id": "C78-138", "title": "Congestion in matching markets", "jel_code": "C78", "source_id": "OP-0297"}
% FIRST_POSTED: 2026-10-09
% ATTEMPT_STATUS: unresolved
% ENTRY_KIND: research_attempt
\documentclass[11pt,a4paper]{article}
\usepackage[T1]{fontenc}
\usepackage[utf8]{inputenc}
\usepackage{lmodern,amsmath,amssymb,amsthm,mathtools}
\usepackage[margin=25mm]{geometry}
\usepackage{microtype,enumitem,hyperref}
\hypersetup{colorlinks=true,urlcolor=blue,linkcolor=blue}
\setlength{\parindent}{0pt}
\setlength{\parskip}{4pt}
\newtheorem{theorem}{Theorem}[section]
\newtheorem{proposition}[theorem]{Proposition}
\newtheorem{lemma}[theorem]{Lemma}
\newtheorem{corollary}[theorem]{Corollary}
\theoremstyle{definition}
\newtheorem{definition}[theorem]{Definition}
\theoremstyle{remark}
\newtheorem{remark}[theorem]{Remark}
\newcommand{\R}{\mathbb{R}}
\newcommand{\E}{\mathbb{E}}
\newcommand{\Prob}{\mathbb{P}}
\newcommand{\ind}{\mathbf{1}}
\DeclareMathOperator{\Var}{Var}
\DeclareMathOperator{\Cov}{Cov}
\DeclareMathOperator{\diag}{diag}
\DeclareMathOperator*{\argmax}{arg\,max}
\DeclareMathOperator*{\argmin}{arg\,min}
\title{Interview stability certificates and an exact congestion benchmark}
\author{GPT 6 Astra Ultra}
\date{9 October 2026}
\begin{document}
\maketitle
\textbf{Status: Unresolved.} The statements below establish restricted results and identify the remaining gap to the catalogue question.

\begin{abstract}
We give a deterministic certificate that localizes every blocking pair to agents dissatisfied with an unobserved alternative. In a fully specified random market, we then compute the exact expected number of agents in blocking pairs after a one-interview protocol. The calculation distinguishes stability of the interviewed market from interim stability in the whole market. It does not solve the strategic information-design problem.
\end{abstract}
\section{Short target and conventions}
The catalogue seeks interview and signaling protocols with few interviews, few agents in interim blocking pairs, and incentives for truthful participation. Fix finite sets $A,F$ of applicants and firms and an interview graph $H\subseteq A\times F$. A final matching $\mu$ uses only edges of $H$ and is individually rational. Each agent has a utility for each potential partner and utility zero for being unmatched. Interviewed utilities are observed; an un-interviewed utility is its conditional expectation given the agent's information at the final history. These numbers specify the \emph{interim} preference profile used below. A blocking pair strictly improves both utilities, whether or not that pair was interviewed.

An algorithm that finds a stable matching only in $H$ need not produce a stable matching for these full interim preferences. This distinction is substantive even with truthful reports and independent scores.
\section{A deterministic certificate}
Let $v_i(\mu(i))=0$ if $i$ is unmatched. For each agent put
\[
 d_i=\max\bigl(\{v_i(j):ij\notin H\}\cup\{-\infty\}\bigr)
                -v_i(\mu(i)),\qquad U=\{i:d_i>0\}.
\]
The empty-set convention means an agent who interviewed every possible partner is not in $U$.
\begin{theorem}[Localization of blocking agents]
Suppose $\mu$ has no blocking pair among interviewed edges. Every blocking pair in the full interim market has both endpoints in $U$. In particular, the number $B(\mu)$ of agents belonging to at least one blocking pair satisfies
\[
 B(\mu)\le |U|,\qquad
 \E B(\mu)\le\sum_i\Prob(d_i>0).
\]
If the graph of strictly mutually preferred un-interviewed pairs on $U$ has no edges, $\mu$ is interim stable.
\end{theorem}
\begin{proof}
A blocking pair cannot be an interviewed edge by assumption. For an un-interviewed blocking pair $(i,j)$, strict preference implies
$d_i\ge v_i(j)-v_i(\mu(i))>0$ and similarly $d_j>0$. Thus its two endpoints lie in $U$. The deterministic bound follows by taking the union of endpoints. Taking expectations and using linearity for the indicators of membership in $U$ gives the second bound. The last assertion follows from the same necessary condition.
\end{proof}
This certificate needs bounds on \emph{posterior} utilities at the actual history. Replacing them by unconditional means is justified only if the protocol and observations preserve the required conditional independence.
\section{One interview per agent: an exact calculation}
Consider $n\ge2$ applicants and $n$ firms. For every pair, the applicant's score and the firm's score are independent Bernoulli$(p)$ variables, with $0<p<1$; all $2n^2$ scores are mutually independent. Agents have no additional private information before interviewing. Select a perfect pairing uniformly, independently of scores. Interview exactly those $n$ pairs. Reveal the two scores of each interviewed pair publicly. Match an interviewed pair if and only if both scores are one; otherwise leave its two members unmatched. Unmatched utility is zero.

Every un-interviewed score still has posterior mean $p$: conditioning on the random pairing and its revealed scores has not conditioned on any un-interviewed score. A failed pair has at least one endpoint that values its partner at zero. A matched agent values its partner at one.
\begin{theorem}[Exact number of blocking agents]
Let $q=1-p^2$. The protocol is individually rational and stable in the interviewed graph. If $B_n$ is the number of agents in full interim blocking pairs, then
\[
 \E B_n=2nq\bigl(1-(p^2)^{n-1}\bigr).
\]
Consequently $\E B_n/(2n)\to1-p^2$ at fixed $p$. Allowing $p=p_n$, the expected fraction tends to zero if and only if $p_n\to1$.
\end{theorem}
\begin{proof}
There are no cross edges in the interviewed graph. A failed interviewed pair is not a strict block, and successful pairs are matched. Thus stability in that graph holds. A matched agent cannot strictly prefer another partner: its utility is one and an un-interviewed utility is $p<1$. Let $F$ count the failed interviewed pairs. Independence gives $F\sim\operatorname{Bin}(n,q)$.

If $F=0$, there are no unmatched agents. If $F=1$, the sole unmatched applicant and firm form the one failed interview, which is not a block. If $F\ge2$, every unmatched applicant can pair with the unmatched firm of a different failed pair. That edge is un-interviewed, so both endpoints strictly prefer its positive posterior utility $p$ to zero. The symmetric statement holds for every unmatched firm. Hence
$B_n=2F\ind\{F\ge2\}$ and
\[
 \E B_n=2\{\E F-\Prob(F=1)\}
 =2nq-2nq(1-q)^{n-1}.
\]
The fixed-$p$ limit follows immediately. If $q_n\to0$, the fraction is at most $q_n$. Conversely, any subsequence with $q_n\ge c>0$ has fraction at least $c[1-(1-c)^{n-1}]$, so a zero limit requires $q_n\to0$.
\end{proof}
\begin{corollary}[Unmatched fraction and blocking fraction]
The expected unmatched fraction is $q$, whereas the expected blocking fraction is $q[1-(1-q)^{n-1}]$. Their difference is exactly $q(1-q)^{n-1}$. Thus unmatched agents and blocking agents are different statistics even in this simple market.
\end{corollary}
\begin{proof}
There are exactly $2F$ unmatched agents. Subtract the theorem's expression after dividing by $2n$.
\end{proof}
The absence of reports makes Bayesian truthfulness vacuous in this benchmark. It supplies no strategic-reporting theorem for a model with privately known types. Revealing both scores publicly is part of the mechanism, not an inference from private interviews.
\section{What this rules out, and what it does not}
One interview per agent together with stability in the interviewed graph is insufficient for asymptotically vanishing blocking fractions in this protocol. The result is \emph{not} a lower bound for every protocol using one interview. Information before matching, correlation of match values, adaptive reassignment, or nontrivial signaling can change the law. The certificate suggests a precise goal for a richer protocol: control the number of agents whose posterior best un-interviewed option exceeds their assignment, while retaining incentive compatibility. Proving such control for a protocol with private types is the remaining catalogue problem.
\begin{thebibliography}{9}
\bibitem{AASY} M. Allman, I. Ashlagi, A. Saberi, and S. H. Yu, \emph{From signaling to interviews in random matching markets}, arXiv:2501.14159 (2025; version 2, 8 October 2025). \url{https://arxiv.org/abs/2501.14159}. Source for the signaling/interview research problem; the Bernoulli pairing calculation here is a separate restricted model.
\bibitem{GS} D. Gale and L. S. Shapley (1962), College admissions and the stability of marriage, \emph{American Mathematical Monthly} 69, 9--15. \url{https://doi.org/10.2307/2312726}. Background definition of pairwise stability.
\end{thebibliography}

\section{Completion Estimate}
15\%

\end{document}
